\documentclass[11pt]{article}
\usepackage[a4paper,margin=1.9cm]{geometry}
\usepackage[T1]{fontenc}
\usepackage{textcomp}
\usepackage[spanish]{babel}
\usepackage{amsmath,amssymb}
\usepackage{lmodern}
\renewcommand{\familydefault}{\sfdefault}
\usepackage{enumitem}

\newcommand{\vect}[2]{\begin{pmatrix}#1\\#2\end{pmatrix}}
\newcommand{\vectiii}[3]{\begin{pmatrix}#1\\#2\\#3\end{pmatrix}}
\usepackage{tikz}
\usetikzlibrary{calc,arrows.meta,patterns,decorations.pathmorphing}
\usepackage[table]{xcolor}
\usepackage{multicol}
\usepackage{parskip}

\definecolor{unalblue}{RGB}{31,73,124}

\newlist{opciones}{enumerate}{1}
\setlist[opciones]{label=(\alph*),itemsep=6pt,topsep=6pt,leftmargin=1.8em,font=\normalfont}
\setlist[enumerate,1]{itemsep=1.4em,topsep=1em}

\newcommand{\examheader}{%
  \noindent
  \begin{minipage}[t]{0.6\linewidth}
    {\small\bfseries Geometría Vectorial y Analítica --- Parcial 3}\\
    {\small Semestre 2014--02}
  \end{minipage}%
  \hfill
  \begin{minipage}[t]{0.35\linewidth}
    \raggedleft\small Escuela de Matemáticas. Universidad Nacional de Colombia.
  \end{minipage}\\[2pt]
  \rule{\linewidth}{0.6pt}
}
\newcommand{\examfooter}{%
  \vfill
  \rule{\linewidth}{0.4pt}\\[1pt]
  {\scriptsize\textit{Transcripción digital --- MathUNAL}\hfill\textcolor{unalblue}{\textbf{\thepage}}}
}
\pagestyle{empty}

\newcommand{\nota}[1]{{\small\itshape\color{unalblue!70!black} #1}}

\begin{document}
\examheader

\begin{center}
Universidad Nacional de Colombia\\
Facultad de Ciencias --- Escuela de Matemáticas\\
\textbf{Tercer examen parcial de Geometría Vectorial}\\
24 de noviembre de 2014
\end{center}

\vspace{10pt}
\textbf{PUNTAJE. SOLO PARA USO OFICIAL}

\begin{center}
\begin{tabular}{|c|c|c|c|c|}
\hline
1 & 2 & 3 & 4 & TOTAL \\
\hline
 & & & & \\
\hline
\end{tabular}
\end{center}

\vspace{10pt}
\textbf{Instrucciones:} No está permitido el uso de calculadora ni de otros aparatos electrónicos. Solucione las preguntas en los espacios en blanco asignados a cada una de ellas. Toda respuesta tiene que estar respaldada por un procedimiento para ser valorada. La última página puede usarse para operaciones en borrador. No está permitido sacar hojas en blanco ni ningún tipo de apuntes durante el examen. \textbf{Duración del examen: 1 hora y 50 minutos.}

\begin{enumerate}
\item (25 puntos) Demuestre que las rectas $\mathcal{L}_1$ y $\mathcal{L}_2$ con ecuaciones:
$$\mathcal{L}_1:\begin{cases}x=2+5t,\\ y=-1+t,\\ z=-3-t,\end{cases} t\in\mathbb{R}\qquad\qquad \mathcal{L}_2:\frac{x-1}{4}=\frac{z-1}{-2},\ y=-\frac{2}{3}$$
se cruzan. (Si necesita más espacio para escribir la solución utilice la última página)

\item Dados los puntos $A=\vect{1}{0}$ y $B=\vect{0}{1}$:

\begin{enumerate}[label=(\alph*)]
\item (10 puntos) Encuentre una ecuación de segundo grado para describir el lugar geométrico de todos los puntos $P=\vect{x}{y}$ tales que el producto de la pendiente de la recta que pasa por $A$ y $P$ multiplicada por la pendiente de la recta que pasa por $B$ y $P$ es constante e igual a $-4$. Identifique a cuál cónica corresponde el lugar geométrico descrito, y encuentre su centro o vértice según corresponda.
\end{enumerate}

Considere las siguientes ecuaciones de segundo grado:
$$\text{(A) } \frac{2(x-1)^2}{16}+\frac{(y+2)^2}{9}=1 \qquad \text{(B) } (y-5)^2=\frac{1}{18}(x+5) \qquad \text{(C) } \frac{(y-3)^2}{10}-\frac{(x+2)^2}{4}=-1$$
Completar:

\begin{enumerate}[label=(\roman*)]
\item (5 puntos) La ecuación (A) representa una \underline{\hspace{2.5cm}} con centro $C=$ \underline{\hspace{1.5cm}} y eje focal paralelo al eje \underline{\hspace{1cm}}.
\item (5 puntos) La ecuación (C) representa una \underline{\hspace{2.5cm}} con centro $C=$ \underline{\hspace{1.5cm}} y excentricidad $e=$ \underline{\hspace{1.5cm}}.
\item (5 puntos) La ecuación (B) representa una \underline{\hspace{2.5cm}} con eje focal paralelo al eje \underline{\hspace{1cm}} y con longitud de lado recto \underline{\hspace{1.5cm}}.
\end{enumerate}

\item Se sabe que las rectas $\mathcal{L}_1$ y $\mathcal{L}_2$ dadas por:
$$\mathcal{L}_1: x=y+1=-\frac{z}{2}\qquad\qquad \mathcal{L}_2:\frac{x+1}{3}=-y,\ z=1$$
se cruzan.

\begin{enumerate}[label=(\alph*)]
\item (15 puntos) Encuentre una ecuación general para el plano que contiene a $\mathcal{L}_1$ y es paralelo a $\mathcal{L}_2$.
\item (10 puntos) Calcule la distancia entre $\mathcal{L}_1$ y $\mathcal{L}_2$.
\end{enumerate}

\item Considere los puntos $A=\vectiii{1}{0}{-1}$, $B=\vectiii{0}{1}{1}$, $C=\vectiii{-1}{1}{0}$ y $D=\vectiii{1}{1}{1}$.

\begin{enumerate}[label=(\alph*)]
\item (15 puntos) Pruebe que los puntos $A$, $B$, $C$ y $D$ no son coplanares (es decir, no están situados en un mismo plano).
\item (10 puntos) Halle una ecuación general del plano determinado por los puntos $A$, $C$ y $D$.
\end{enumerate}

\end{enumerate}

\examfooter
\end{document}
